% =====================================================================
%  Annexe A : Installation pas a pas
% =====================================================================
\chapter{Installation pas à pas}
\label{ann:installation}

\chapitreresume{Cette annexe reprend l'installation complète, du téléchargement de
Python à la première page du projet, avec les commandes pour Windows, macOS et
Linux, ainsi que les solutions aux problèmes les plus courants.}

\section{Installer Python}

\subsection{Windows}

\begin{enumerate}
  \item Téléchargez Python 3.12 sur \url{https://www.python.org/downloads/}.
  \item Lancez l'installateur.
  \item Cochez \og Add python.exe to PATH \fg{} avant de cliquer sur
        \og Install Now \fg.
  \item À la fin, ouvrez PowerShell et vérifiez :
\end{enumerate}

\begin{terminal}
python --version
python -m pip --version
\end{terminal}

\begin{attention}
Sous Windows, \code{python} peut ouvrir le Microsoft Store si Python n'est pas
installé. Tapez \code{py --version} pour vérifier. Si la commande \code{py}
fonctionne mais pas \code{python}, utilisez \code{py} partout dans ce module, ou
réinstallez Python en cochant l'option PATH.
\end{attention}

\subsection{macOS}

\begin{terminal}
brew install python@3.12
python3 --version
\end{terminal}

\begin{note}
macOS fournit parfois une version ancienne de Python. Utilisez \code{python3} et
\code{pip3} pour éviter les confusions. Homebrew installe une version à jour.
\end{note}

\subsection{Linux (Debian, Ubuntu)}

\begin{terminal}
sudo apt update
sudo apt install python3 python3-venv python3-pip
python3 --version
\end{terminal}

Sur Fedora ou Red Hat, remplacez \code{apt} par \code{dnf}. Sur Arch,
\code{pacman -S python}.

\section{Créer l'environnement virtuel}

Un environnement virtuel isole les paquets du projet. Chaque projet peut avoir sa
propre version de Django sans conflit.

\begin{terminal}
mkdir gesibiblio
cd gesibiblio
python -m venv .venv
\end{terminal}

Activez-le :

\begin{terminal}
# Windows PowerShell
.\.venv\Scripts\Activate.ps1

# Windows cmd
.venv\Scripts\activate.bat

# macOS et Linux
source .venv/bin/activate
\end{terminal}

Le prompt affiche \code{(.venv)} au début. C'est le signe que l'environnement est
actif.

\begin{attention}
Sous PowerShell, l'activation peut être bloquée par la politique d'exécution. Dans
ce cas, ouvrez PowerShell en administrateur et lancez :
\code{Set-ExecutionPolicy -Scope CurrentUser RemoteSigned}. Fermez puis rouvrez le
terminal.
\end{attention}

\begin{bonne}
Ne versionnez jamais le dossier \code{.venv/}. Ajoutez-le à \code{.gitignore}. Un
nouveau développeur recrée son environnement avec \code{requirements.txt}.
\end{bonne}

\section{Installer Django et Pillow}

\begin{terminal}
python -m pip install --upgrade pip
python -m pip install "Django==5.2" Pillow
python -m django --version
\end{terminal}

Pillow est nécessaire pour \code{ImageField}. Sans lui, Django refuse de démarrer
dès qu'un modèle contient une image.

\begin{terminal}
5.2
\end{terminal}

\begin{note}
Les guillemets autour de \code{"Django==5.2"} évitent que le shell interprète les
signes égal. Sous PowerShell, les guillemets sont recommandés.
\end{note}

\section{Créer le projet}

\begin{terminal}
django-admin startproject gesibiblio .
\end{terminal}

Le point final est important : il crée les fichiers dans le dossier courant au lieu
de créer un sous-dossier \code{gesibiblio/gesibiblio/}.

Vérifiez l'arborescence :

\begin{codebox}[tree]{Après startproject}
.
|-- manage.py
`-- gesibiblio/
    |-- __init__.py
    |-- settings.py
    |-- urls.py
    |-- asgi.py
    `-- wsgi.py
\end{codebox}

Lancez le serveur :

\begin{terminal}
python manage.py runserver
\end{terminal}

Ouvrez \code{http://127.0.0.1:8000/}. La fusée de Django s'affiche. Le projet
fonctionne.

\section{Les premiers réglages}

\begin{codebox}[python]{gesibiblio/settings.py, extraits}
LANGUAGE_CODE = "fr-fr"
TIME_ZONE = "Africa/Douala"
USE_I18N = True
USE_TZ = True
\end{codebox}

\begin{description}[leftmargin=4.3cm,style=nextline]
  \item[\code{LANGUAGE\_CODE}] Langue de l'interface d'administration et des
  messages.
  \item[\code{TIME\_ZONE}] Fuseau horaire. Adaptez-le à votre public.
  \item[\code{USE\_I18N}] Active la traduction.
  \item[\code{USE\_TZ}] Stocke les dates en UTC et les convertit à l'affichage.
\end{description}

\begin{bonne}
Laissez \code{USE\_TZ = True}. Les dates sont stockées en UTC, ce qui évite les
ambiguïtés lors des changements d'heure. Utilisez \code{timezone.localdate()} dans
le code plutôt que \code{datetime.date.today()}.
\end{bonne}

\section{Initialiser la base}

\begin{terminal}
python manage.py migrate
\end{terminal}

Cette commande crée les tables des applications intégrées : utilisateurs, sessions,
permissions, administration. Elle crée aussi le fichier \code{db.sqlite3}.

\begin{terminal}
python manage.py createsuperuser
\end{terminal}

\section{Configurer Visual Studio Code}

Extensions recommandées :
\begin{itemize}
  \item Python (Microsoft) ;
  \item Pylance ;
  \item Django (Baptiste Darthenay) pour la coloration des gabarits ;
  \item SQLite Viewer pour inspecter la base ;
  \item Even Better TOML, utile pour la configuration.
\end{itemize}

Sélectionnez l'interpréteur de l'environnement virtuel : \code{Ctrl+Shift+P}, puis
\og Python: Select Interpreter \fg, et choisissez \code{.venv}.

\begin{codebox}[tree]{.vscode/settings.json}
{
  "python.defaultInterpreterPath": ".venv/Scripts/python.exe",
  "editor.formatOnSave": true,
  "files.exclude": {
    "**/__pycache__": true
  }
}
\end{codebox}

Sous macOS et Linux, le chemin est \code{.venv/bin/python}.

\section{Installer Git}

\begin{terminal}
git init
git add .
git commit -m "Projet Django initial"
\end{terminal}

Créez un dépôt distant sur GitHub ou GitLab, puis :

\begin{terminal}
git remote add origin https://github.com/utilisateur/gesibiblio.git
git push -u origin main
\end{terminal}

\begin{attention}
Avant le premier commit, vérifiez que \code{.gitignore} exclut \code{.venv/},
\code{db.sqlite3}, \code{.env} et \code{media/}. Un secret poussé sur GitHub reste
dans l'historique, même après suppression du fichier.
\end{attention}

\section{Problèmes fréquents}

\begin{table}[htbp]
\centering
\small
\begin{tabularx}{\textwidth}{L{5.1cm} X}
\toprule
\textbf{Message} & \textbf{Solution} \\
\midrule
\code{'python' n'est pas reconnu} & Python n'est pas dans le PATH. Réinstallez en
cochant l'option, ou utilisez \code{py}. \\
\code{No module named django} & L'environnement virtuel n'est pas activé, ou Django
n'est pas installé. Activez \code{.venv} et relancez \code{pip install}. \\
\code{Activate.ps1 cannot be loaded} & Politique PowerShell. Exécutez
\code{Set-ExecutionPolicy -Scope CurrentUser RemoteSigned}. \\
\code{Port 8000 is already in use} & Un autre serveur tourne. Lancez
\code{runserver 8001} ou arrêtez le processus. \\
\code{ModuleNotFoundError: catalogue} & L'application n'est pas dans
\code{INSTALLED\_APPS}, ou \code{startapp} n'a pas été lancé au bon endroit. \\
\code{TemplateDoesNotExist} & Chemin de gabarit incorrect, ou
\code{TEMPLATES} mal configuré. \\
\code{NoReverseMatch} & Nom de route inconnu, ou paramètre manquant. \\
\code{OperationalError: no such table} & Les migrations n'ont pas été appliquées.
Lancez \code{migrate}. \\
\code{CSRF verification failed} & \code{\{\% csrf\_token \%\}} manquant dans le
formulaire. \\
\code{DisallowedHost} & Ajoutez le domaine à \code{ALLOWED\_HOSTS}. \\
\bottomrule
\end{tabularx}
\caption{Erreurs d'installation et solutions.}
\end{table}

\begin{piege}
Beaucoup d'erreurs viennent d'un environnement virtuel oublié. Avant de chercher
une cause compliquée, vérifiez que le prompt contient \code{(.venv)} et que
\code{which python} (ou \code{where python} sous Windows) pointe vers le dossier du
projet.
\end{piege}

\section{Récapitulatif des commandes}

\begin{codebox}[bash]{Cycle de travail quotidien}
# Activer l'environnement
.\.venv\Scripts\Activate.ps1     # Windows
source .venv/bin/activate         # macOS et Linux

# Lancer le serveur
python manage.py runserver

# Après une modification de modèle
python manage.py makemigrations
python manage.py migrate

# Vérifier le projet
python manage.py check
python manage.py test

# Créer un compte administrateur
python manage.py createsuperuser

# Ouvrir un shell Python connecté à Django
python manage.py shell
\end{codebox}
